\documentclass[11pt,a4paper]{article}
\usepackage[utf8]{vietnam}
\usepackage{amsmath,amssymb,amsthm}
\usepackage{geometry}
\geometry{a4paper, margin=2.5cm}

\theoremstyle{definition}
\newtheorem{theorem}{Định lý}
\newtheorem{definition}{Định nghĩa}

\title{\textbf{Ghi chú bài giảng Giải tích \& Đại số}}
\author{}
\date{}

\begin{document}

\maketitle

\section{Định lý kẹp (Sandwich Theorem)}

\begin{theorem}
Cho ba hàm số $f(x)$, $g(x)$, $h(x)$ thỏa mãn:
\[
f(x) \leq g(x) \leq h(x)
\]
với mọi $x$ thuộc một lân cận của điểm $a$. 

Nếu:
\[
\lim_{x \to a} f(x) = \lim_{x \to a} h(x) = L
\]
thì:
\[
\lim_{x \to a} g(x) = L.
\]
\end{theorem}

\section{Khai triển Taylor -- Maclaurin cơ bản}

Khi $x \to 0$, ta có các khai triển Maclaurin chuẩn sau:

\begin{align*}
e^x &= 1 + x + \frac{x^2}{2!} + \frac{x^3}{3!} + \dots + \frac{x^n}{n!} + o(x^n) \\
\sin(x) &= x - \frac{x^3}{3!} + \frac{x^5}{5!} - \dots + (-1)^n \frac{x^{2n+1}}{(2n+1)!} + o(x^{2n+2})
\end{align*}

\section{Bài tập về nhà}

\begin{enumerate}
    \item[\textbf{Câu 1:}] Tính giới hạn:
    \[
    I = \lim_{x \to 0} \frac{e^x - 1 - x}{x^2}
    \]

    \item[\textbf{Câu 2:}] Cho ma trận:
    \[
    A = \begin{pmatrix}
    1 & 2 & 3 \\
    0 & 1 & 4 \\
    5 & 6 & 0
    \end{pmatrix}
    \]
    Tìm định thức $\det(A)$ và ma trận nghịch đảo $A^{-1}$.
\end{enumerate}

\end{document}